\section{De Base Model a Chatbot: tres señales de supervisión}

La sección anterior estableció los límites de la evidencia; esta sección construye primero un lenguaje común. Lo que se llama «entrenar un chatbot» comprende al menos tres niveles: demonstration learning, preference learning y policy optimization. Cada uno usa estructuras de datos distintas, resuelve problemas distintos y requiere también métodos de evaluación distintos.

\subsection{SFT, Reward Model y RLHF}

\term{SFT (supervised fine-tuning, fine-tuning supervisado)} toma el prompt $x$ y una completion $y$, humana o sintética, como par supervisado, y minimiza la next-token cross-entropy:
\[
\mathcal{L}_{\mathrm{SFT}}(\theta)
=-\sum_{t=1}^{|y|}\log \pi_{\theta}(y_t\mid x,y_{<t}).
\]
Aquí $\theta$ son los parámetros del chatbot, $y_t$ es el $t$-ésimo token objetivo y $y_{<t}$ es el prefijo ya dado. Lo que aprende SFT es «responder como en la demostración»; no aprende directamente a comparar cuál de dos respuestas es mejor.

\lecturefigure{slide-04.jpg}{Flujo de entrenamiento al estilo InstructGPT: SFT proporciona el punto de partida de la helpfulness}{Diapositivas oficiales p.4}

\begin{knowledgebox}{Lectura de la figura: la estructura de las label es distinta en cada paso}
En el primer paso, personas escriben el prompt y la demonstration; en el segundo, el modelo genera varias respuestas y las personas solo las ordenan; en el tercero, la learned reward actualiza la policy. Primero observe quién produce el texto en cada paso y después cómo el trabajo humano pasa de «escribir la respuesta completa» a «comparar respuestas»; así se entiende por qué la escala, el ruido y el costo de los datos difieren entre las tres etapas.
\end{knowledgebox}

\term{reward model (modelo de recompensa, RM)} asigna a $(x,y)$ un escalar $r_{\phi}(x,y)$. Para una preferred answer $y_w$ y una rejected answer $y_l$, suele usarse la pairwise loss de Bradley--Terry:
\[
\mathcal{L}_{\mathrm{RM}}(\phi)
=-\log \sigma\!\left(r_{\phi}(x,y_w)-r_{\phi}(x,y_l)\right).
\]
Aquí $\phi$ son los parámetros del reward model y $\sigma$ es la sigmoid; al modelo le basta con que la puntuación del winner supere a la del loser, sin que el valor de la recompensa tenga un significado físico absoluto.

\lecturefigure{slide-05.jpg}{Preference ranking y RLHF añaden harmlessness y restricciones de valores}{Diapositivas oficiales p.5}

\begin{importantbox}{El objetivo de RLHF no es «maximizar la reward y listo»}
\term{RLHF (reinforcement learning from human feedback)} optimiza la policy con la learned reward y, por lo general, la restringe al mismo tiempo con una KL penalty para que no se aleje del reference model:
\[
\max_{\theta}\;\mathbb{E}_{y\sim\pi_{\theta}(\cdot\mid x)}
\left[r_{\phi}(x,y)-\beta\log\frac{\pi_{\theta}(y\mid x)}{\pi_{\mathrm{ref}}(y\mid x)}\right].
\]
Aquí $\beta$ controla la tensión entre «perseguir la preferencia» y «conservar el comportamiento del modelo original». El reward hacking, el distribution shift y los errores del RM pueden hacer que una puntuación más alta no equivalga a una mejor respuesta.
\end{importantbox}

\teachervoice{El ponente explica el step 1 como «hacer primero que el modelo sea un helpful chatbot» y presenta los steps 2/3 como la incorporación de guardrails. Esta formulación es un relato de clase; no significa que SFT solo afecte la helpfulness ni que RLHF solo afecte la harmlessness: en la práctica, las tres etapas modifican a la vez el estilo, la exactitud factual y el comportamiento de rechazo.}

\subsection{Términos clave: no confundir los tres tipos de muestras}

La distinción siguiente es el conocimiento previo mínimo para leer las figuras posteriores. Si a la demonstration, al preference pair y al policy rollout se les llama por igual «datos de entrenamiento», no se puede explicar por qué la misma cantidad de muestras tiene costos y contenido informativo completamente distintos.

\begin{knowledgebox}{Tres tipos de señales de supervisión}
\begin{tabularx}{\textwidth}{@{}L{0.18\textwidth}L{0.25\textwidth}X@{}}
\toprule
Objeto & Estructura de la muestra & Qué aprende \\
\midrule
SFT & prompt + target completion & Imitar la token distribution de la respuesta objetivo \\
Reward model & prompt + chosen/rejected responses & Aprender un orden de preferencia relativo \\
RL policy & prompt + sampled response + scalar reward & Hacer que la policy produzca con más frecuencia responses de reward alta \\
\bottomrule
\end{tabularx}
\end{knowledgebox}

\subsection{Resumen de la sección}

SFT, reward modeling y RLHF usan tres tipos distintos de label. Cuando más adelante se hable de 10K demonstrations, 20K dialogues o 100K preference examples, primero hay que preguntar a qué estructura de muestra pertenecen; no se pueden comparar directamente las cifras.
